\documentclass[12pt,a4paper]{article}

% ================================================================
% PAGE SETUP 
% ================================================================

\usepackage[a4paper,left=2cm,right=2cm,top=2cm,bottom=2cm]{geometry}
\usepackage{setspace}
\usepackage{changepage}
\usepackage{titlesec}

\setcounter{secnumdepth}{0}
\titleformat{\section}
  {\normalfont\Large\bfseries\centering}
  {}{0pt}{}

% ================================================================
% TYPESETTING AND LANGUAGE
% ================================================================

\usepackage[T1]{fontenc}
\usepackage{newtxtext,newtxmath}
\usepackage[ngerman]{babel}
\usepackage{csquotes}
\usepackage[normalem]{ulem}
\newenvironment{longquote}
{%
    \vspace{-1.5em}
    \begin{adjustwidth}{2cm}{2cm}
        \fontsize{10}{12}\selectfont
        \singlespacing
        \itshape
    }
{%
    \end{adjustwidth}
}

% ================================================================
% FOOTNOTES AND BIBLIOGRAPHY   
% ================================================================

\usepackage[
    backend=biber,
    style=deutschstellenbosch,
    autocite=footnote
]{biblatex}

\addbibresource{quellen.bib}

\makeatletter
\renewcommand{\footnotesize}{%
    \def\baselinestretch{1}%
    \@setfontsize\footnotesize{10}{12}%
}
\makeatother

\usepackage[hidelinks,hyperfootnotes=true]{hyperref}

\begin{document}
\onehalfspacing

% ================================================================
% BODY
% ================================================================

\noindent
\textbf{%
    Quin Barnard\\
    27902277\\
    Deutsch 358: Gliederung für die Abschlussarbeit\\
    Herr Professor Ingo Müller
}

\par
\section{Einleitung}
In der Einleitung wird erwähnt, dass Kafkas Erzählungen bereits intensiv anhand von Affekten wie Angst, Scham und Schuld analysiert worden sind. Die beiden Kurzgeschichten \textit{„Kinder auf der Landstraße“} und \textit{„Unglücklichsein“}, die im Band \textit{„Betrachtung“} enthalten sind, beinhalten sowohl diese als auch weniger drastische und weniger bemerkte Affekte. Die Protagonisten beider Geschichten erwarten scheinbar ihren jeweiligen Gegner, jedoch kommen diese zuerst unbekannt vor und streiten sich mit den Protagonisten. Es wird hier eine gewisse Inkongruenz zwischen was der Leser mitbekommt, was der Protagonist offensichtlich erfährt und was die Gegner erfahren thematisiert. Hier wird gezeigt, wie der Mensch „über die Wirklickheit keine wahren, sondern nur noch wahrscheinliche Aussagen“\autocite[137]{horstkotte2014} äußern kann. Das epistemologische Problem ist selbstverständlich mit Furcht und Angst zu verbinden, aber ebenso interessant ist wie das Problem durch die Affekte wie Unbehagen, Unsicherheit, Verlegenheit, Unentschlossenheit usw. sich gestalten lässt. In beiden Geschichten ist die Erlebnis des lyrischen Ichs unmittelbar hierdurch geprägt, obwohl es dem Leser nicht unbedingt bei der ersten Lektüre auffällt. Die anderen Texte im Band \textit{„Betrachtung“} stellen ähnliche Situationen dar, \textit{„Der Fahrgast“} z.B. versteht nicht warum das Mädchen sich nicht ebenso unwohl fühlt wie er. Daher könnte bei der Interpretation der zwei längsten Texte die Gesamtheit herbeigezogen wird für eine kohärente Interpretation. Es wird aber häufig auf \textit{„Kinder auf der Landstraße“} und \textit{„Unglücklichsein“} fokusiert, da ihre Längen eine detaillierte Auseinandersetzung erlauben.\par Die Forschungsfrage(n) werden gestellt:\\ \uline{Wie überträgt die Erzählweise die Unsicherheit der Figur auf den Leser?}\\ und\\ \uline{Wie wird diese Verunsicherung der Figur von Affekten wie Zweifel, Verlegenheit und Unbehagen geprägt?}

\section{Teil 1}
Der Aufsatz fängt mit einer Diskussion der Titeln an. Elizabeth Boa\autocite{boa2014}, unter anderen\autocite[see 105]{fromm}, bietet eine Interpretation vom Namen \textit{„Betrachtung“} die auf dessen Zweideutigkeit, nämlich „Betrachtung“ als hinausschauen und „Betrachtung“ als reflexives Denken, hinweist. \textit{„Kinder auf der Landstraße“} ist einfach zu deuten, weil die Erzählung tatsächlich nichts als Kinder auf einer Landstraße darstellt. \textit{„Unglücklichsein“} ist aber sehr abstrakt und erst schwierig zu interpretieren.\par 

\vspace{0.5em}
\begin{longquote}
»Kafka gibt dem Erzähler-Ich die Propleme des Autors zur Lösung auf. Ein Problem des Textes \textit{Unglücklichsein} ist es, ein Paradox zu lösen, das nicht lösbar erscheint. So sagt das Erzähler-Ich zum Kind, sie seien von der gleichen Natur. Das Paradox besteht darin, Gespenst und zugleich Subjekt zu sein. Die Schrift verhält sich zum Leben wie das Gespenst zum Subjekt. Will man in Kafkas Bildlichkeit verbleiben, ließe sich paradox formulieren: Das Gespenst der Schrift ist wirklich unwirklich«\autocite[117]{fromm}
\end{longquote}

Fromms Perspektive bringt sofort das epistemologische Problem hervor und kann auch der Titel als Bezeichnung für die dementsprechende Stimmung erklären.

\section{Teil 2}
Zunächst wird \textit{„Kinder auf der Landstraße“} genauer untersucht. Interessant ist, dass das Ich seufzt als die laufenden Kinder erscheinen, obwohl der Leser keinen Grund für das Seufzen mitbekommt. Eines der Kinder diese Frage, aber ist auch gleich dem Ich gegenüber feindselig. Der Wechsel der Perspekive, wie z.B. \textit{„[…]bis ich nicht mehr glaubte, daß sie stiegen, sondern daß ich falle“}\footnote{Kafkas Band \textit{Betrachtung} wird hier und im Folgenden zitiert nach: \cite{kafka} Hier S. 6.} und die Identität als Mitgleid einer Gruppe werden untersucht. Der Zweifel, der oft in dieser Geschichte auftaucht, wird nach Nietzsches bekanntes Zitat \textit{„Es gibt keine Fakten, nur Interpretationen“} analysiert, wie Gosetti-Ferencei\autocite[22]{gosetti2021} und Horstkotte\autocite[137]{horstkotte2014} es auch getan haben. Es wird auch besprochen, wie eines der laufenden Kinder ständig dem lyrischen Ich Vorwürfe, wie \textit{„Was versteckst Du Dich, laß den Unsinn!“}(8), macht. Es scheint dem Leser kaum einen Grund dafür zu geben, doch ist die Erlebnis des Ich davon geprägt. Zunächst wird die Kausalität an mancher Stellen hinterfragt. Es heißt im Text: \textit{„Das Wasser unten schlug an Steine und Wurzeln, als wäre es nicht schon spät Abend.“} Freilich gibt es keinen Grund, warum das Wasser nicht im spät Abend schlagen darf. Zuletzt wird das Ende der Erzählung besprochen nach Fromms\autocite[113]{fromm} Interpretation: die Narren seien außerhalb der Wirklichkeit und stehen in einem mythischen Bereich.

\section{Teil 3}
Die Kausalität ist auch zu hinterfragen bei \textit{„Unglücklichsein“}. Fromm stellt die Frage: \textit{„[Warum beruhigt] der Blick zum Fenster[?]“}\autocite[111]{fromm} Der Text fängt sogar mit einem kaum verständlichen Satz an, der den Leser hinterfragen lässt, was überhaubt unerträglich geworden ist. Der Erzähler schreit, nur um etwas zu hören, was an Horstkottes \textit{„[erwünschtem Verschwinden] des Subjekts in der eigenen Arbeit“}\autocite[122]{horstkotte2014} erinnert. Es scheint als ob das Ich sich unsicher fühlt: einerseits in gewissermaßen physisch unsicher, da er behauptet die Eile sei nötig(32), andererseits psychologisch unsicher und an seinen Gegenständen zweifelnd, da er sich durch das Schreien beruhigen muss. Das Kind, das erscheint, macht dem lyrischen Ich hier auch Vorwürfe, wie: \textit{„Sie [haben] mir schon gedroht}(34), die gar nicht verständlich sind. Das Ich und das Kind erfahren unterschiedliche „Realitäten“, die der Leser nicht sehen kann. Deleuze und Guattaris Behauptung, dass Kafka nur Experimente protokolliere und Erfahrungen berichte, ohne sie zu deuten\autocite[12]{deleuze2012} wird hier herbeigezogen.

\subsection{Teil 4}
\begin{longquote}
»In der ersten Gesichete, \textit{Kinder auf der Landstraße}, heißt es: „wir bekamen in die Beine einen Galopp wie niemals, bei den Sprüngen hob uns in den Hüften der Wind“. In der letzten Geschichte \textit{Unglücklichsein} heißt es: „Als es schon unerträglich geworden war […] und ich über den schmalen Teppich meines Zimmers wie in einer Rennbahn einherlief“. Kindlich-naive Imagination im Spiel und die Imaginationen eines sozial isolierten Ichs, die ins Schreiben münden, stehen sich gegenüber. Die „wahre Gestalt“ des Ichs entwickelt sich im Umgang mit der Fremdheit seiner Umgebung.«\autocite[103]{fromm}
\end{longquote}

Eine mögliche Verbindung zwischen der Geschichten wird besprochen, um eine Gesamtinterpretation zu bieten. Es ist möglich, dass die zwei Subjekte gleich sind, und es in \textit{„Unglücklichsein“} dem lyrischen Ich unerträglich geworden ist, dass er kein Kind mehr ist. Wäre das der Fall, ist der Punkt im Text, wo das Ich dem Kind sagt, er wisse nicht, wie es später sein werde, besonders interessant, weil er vermutlich sein Wissen zu verstecken versucht. Der Leser kann aber auch nicht wissen, inwiefern die Erlebnisse real oder fantasiert sind. Deswegen ist der Begriff der \textit{„Kindlich-naive Imagination im Spiel“} sehr nützlich beim deuten der epistemologische Frage.

\clearpage
\printbibliography[
    title={Literaturverzeichnis}
]

\end{document} 
